\scanpage{161}\mspage{19.}
% archive sheet number: 161
Par exemple la catégorie $\Delta$, ou tte \inserted{petite} catégorie stable par produits binaires, est tot. $0$-connexe, mais $\Delta$ n'est pas loc. \deleted{\unsure{t.}} irréductible (les catégories induites $\Delta_{/\Delta_n}$, pour $n>0$, ne sont \inserted{\uline{visiblement}} \uline{pas} irréductibles). (Mais rappelons que \deleted{pour} si $\unsure{\mathscr X|\mathcal X}=\mathrm{Top}(T)$, $T$ un espace \inserted{topologique}, alors $T$ irréductible implique $\mathscr X$ \uline{loc.} irréductible. Donc $T$ \unsure{tot.} $0$-connexe implique $\mathrm{Top}(T)$ totalement \uline{et} localement irréductible, i.e. \deleted{\unsure{tout}}\inserted{pour} $T'$ étale sur $T$, tt $0$-connexe, $T'$ est irréductible.)

\inserted{Théorème}\deleted{Proposition} 2.6. Soit $\mathscr X$ un topos \inserted{loc. connexe}, $S\subset\mathscr X^{\wedge}$ une petite partie génératrice \inserted{formée d'objets $0$-connexes$^{*}$} \unsure{.} Conditions équivalentes [\uuline{Sauf} en c'), qui \inserted{\uuline{\unsure{sont|sauf}}} \unsure{semble} plus faible, cf démonstr. bas]

% slanted
\margin{\deleted{\illegible{3}}}

\textcircled{a} Tout topos induit $\mathscr X_{/U}$ \deleted{est tot.} sur $U$ $0$-connexe \inserted{de $\mathscr X^{\wedge}$} est totalement \deleted{\unsure{conn.}} $0$-connexe.

\textcircled{a'} \deleted{\unsure{Le}}\inserted{Tout} produit fibré $U\times_W V$ d'objets $0$-connexes de $\mathscr X^{\wedge}$ est $0$-connexe.

\textcircled{b} $\mathscr X$ est loc. irréductible, et \deleted{totalement}\inserted{ses} \unsure{composantes} connexes sont tot. $0$-connexes

\uline{c)} Comme a), avec $U$ dans $S$. \inserted{Les topos induits $\mathscr X_{/U}$ ($U\in\mathrm{Ob}\,S$) sont totalement $0$-connexes}

\inserted{c') Comme a'), avec $U,V,W$ dans $S$}

% slanted
\margin{\uline{NB} On a l'équivalence de \textcircled{\scriptsize $a,a',b,d$}, qui impliquent c), c'), \deleted{lesquelles} \unsure{égal.} Mais il semble bien que \deleted{ni} c') soit plus faible, cf \unsure{cependant|la question} \unsure{variantes} \textcircled{e} \textcircled{e'} de \unsure{2.6|2.4} (v. \unsure{22, 29|22, 24}).}

% slanted
\margin{C'est bien vrai c) et c') n'impliquent \uuline{pas} la condition \unsure{plus forte} (cf p. 28)}

\textcircled{d} $S$ est loc. filtrante à gauche, i.e. satisfait les deux conditions

\uline{$\mathrm{PS}*1_{\mathrm g}$} Tout diagramme
\diagram{\begin{tikzcd}[ampersand replacement=\&,row sep=small,column sep=small]
a\arrow[dr]\&\\
\&x\\
b\arrow[ur]\&
\end{tikzcd}}
de $S$ s'insère dans
\diagram{\begin{tikzcd}[ampersand replacement=\&,row sep=small,column sep=small]
\&a\arrow[dr]\&\\
c\arrow[ur,dashed]\arrow[dr,dashed]\&\&x\\
\&b\arrow[ur]\&
\end{tikzcd}}
commutatif

\uline{$\mathrm{PS}*2$} Tout diagramme
\diagram{\begin{tikzcd}[ampersand replacement=\&]
a\arrow[r,shift left]\arrow[r,shift right]\& b
\end{tikzcd}}
s'insère dans diagramme commutatif
\diagram{\begin{tikzcd}[ampersand replacement=\&]
c\arrow[r]\&a\arrow[r,shift left]\arrow[r,shift right]\&b
\end{tikzcd}}
